% Source: Wald GR, physical PDF pp. 256--257
%         (printed pp. 243--244).
% Coverage: Chapter 10 opening material before Section 10.1.
% Figures/tables/footnotes: no figures, tables, or footnotes.
% Status: source-reviewed and compile-clean.
% Uncertainties: none; a printed-source typo is corrected below.

% Wald numbers results within each subsection (e.g., Theorem 10.1.1).
\begingroup
\renewcommand{\thetheorem}{\thesubsection.\arabic{theorem}}

As discussed in chapter~4, general relativity asserts that spacetime structure and
gravitation are described by a spacetime \((M,g_{ab})\), where \(M\) is a four-dimensional
manifold and \(g_{ab}\) is a Lorentz metric satisfying Einstein's equation. In chapters~5
and~6 we obtained exact solutions of Einstein's equation which made highly successful
physical predictions concerning cosmology and the structure and gravitation fields of
spherical bodies. However, much more is required in order that general relativity be
a physically viable theory. We see a wide variety of physical phenomena for which
general relativity must account. Thus, it is essential that there exist a correspondingly
wide class of solutions of Einstein's equation. For example, we know on physical
grounds that there is a large class of possible gravitational fields of isolated bodies.
If a correspondingly large class of solutions of Einstein's equation failed to exist, we
would be forced to reject general relativity as a correct theory of nature.

A closely related issue concerns the fact that quite generally in classical physics
we have a great deal of physical control over initial conditions of systems. If the
system then is allowed to evolve freely, its behavior is completely determined by the
initial conditions. For example, in ordinary particle mechanics we are able to control
the initial positions and velocities of the particles. Given these initial conditions, if
the system is permitted to evolve without outside interference, the dynamical evolution
of the particles is determined. Similarly, as discussed in section~10.2, in
electromagnetism we are able to arrange the initial values of the electric field,
\(\mathbf{E}\), and the magnetic field, \(\mathbf{B}\), subject only to the constraints
\(\boldsymbol{\nabla}\mathbin{\cdot}\mathbf{E}=4\pi\rho\) and
\(\boldsymbol{\nabla}\mathbin{\cdot}\mathbf{B}=0\). Again, given these initial conditions, the
subsequent evolution of the system is determined. Although our practical ability to
control initial conditions in gravitational problems is far more limited, it seems natural
to believe that (at least over regions much smaller than cosmological scales) we should,
in principle, be able to control the initial conditions of the gravitational field and matter
distribution, perhaps subject to some constraints as in electromagnetism. Thus, unless
general relativity differs drastically from other theories of classical physics, it should
permit a physically reasonable specification of initial data. Furthermore, given these
initial data, Einstein's equation (possibly supplemented by additional equations for the
matter) should determine the subsequent evolution.

If a theory can be formulated so that ``appropriate initial data'' may be specified
(possibly subject to constraints) such that the subsequent dynamical evolution of the
system is uniquely determined, we say that the theory possesses an initial value
formulation. % SOURCE ERRATUM: printed ``contraints.''
However, even if such a formulation exists, there remain further properties that a
physically viable theory should satisfy. First, in an appropriate sense, ``small changes''
in initial data should produce only correspondingly ``small changes'' in the solution
over any fixed compact region of spacetime. If this property were not satisfied, the
theory would lose essentially all predictive power, since initial conditions can be
measured only to finite accuracy. It is generally assumed that the pathological behavior
which would result from the failure of this property does not occur in physics. Second,
changes in the initial data in a region, \(S\), of the initial data surface should not
produce any changes in the solution outside the causal future, \(J^+(S)\), of this region.
If such changes occurred, we should be able to use them to propagate signals ``faster
than the speed of light.'' This would undermine the entire framework of relativity theory.
If a theory possesses an initial value formulation which satisfies both of the above
properties, we say that this initial value formulation is \emph{well posed}. Note, however,
that we have not attempted to give a mathematically precise definition of ``well posed
initial value formulation'' here since the precise criteria depend on the type of theory
considered.

The purpose of this chapter is to establish that general relativity admits a well posed
initial value formulation. Thus, general relativity survives this rather stringent test of
the physical viability of the theory. We begin in section~10.1 by discussing the initial
value formulation of particle mechanics and the theory of the Klein--Gordon field in
Minkowski spacetime. In section~10.2 we analyze the initial value formulation of
Maxwell's equations in Minkowski spacetime, which shares many features analogous
to general relativity with regard to initial constraints and gauge freedom. The initial
value formulation of general relativity then is presented.

Our analysis of Einstein's equation and other field equations will be based on
expressing them as second order hyperbolic systems. Analysis of Einstein's equation
formulated as a first order hyperbolic system (Fischer and Marsden 1972) will not be
discussed here. In addition, we shall consider only initial value formulations on
spacelike Cauchy surfaces. An initial value formulation for general relativity on an
initial surface formed by two intersecting null surfaces also can be given (Sachs
1962a; Muller zum Hagen and Seifert 1977).
